
\documentclass{article}
\usepackage{colortbl}
\usepackage{makecell}
\usepackage{multirow}
\usepackage{supertabular}

\begin{document}

\newcounter{utterance}

\centering \large Interaction Transcript for game `eqbench', experiment `eq\_bench', episode 133 with reflexion{-}Qwen3.5{-}4B.
\vspace{24pt}

{ \footnotesize  \setcounter{utterance}{1}
\setlength{\tabcolsep}{0pt}
\begin{supertabular}{c@{$\;$}|p{.15\linewidth}@{}p{.15\linewidth}p{.15\linewidth}p{.15\linewidth}p{.15\linewidth}p{.15\linewidth}}
    \# & \multicolumn{2}{c}{Player} && \multicolumn{2}{c}{Game Master} \\
    \hline

    \theutterance \stepcounter{utterance}  
    & & & \multicolumn{4}{p{0.6\linewidth}}{
        \cellcolor[rgb]{0.9,0.9,0.9}{
            \makecell[{{p{\linewidth}}}]{
                \texttt{\tiny{[P1$\langle$GM]}}
                \texttt{You have attempted this same task before and did not succeed. These are the plans you wrote after those attempts.} \\
\texttt{Trial \#1: New plan: I failed because I misinterpreted the emotional arc of the dialogue. Arthur repeatedly apologizes, admits fault, and offers to make amends, which strongly indicates **Remorseful** rather than just Defeated. While he feels defeated by the situation, his primary active emotion driving his actions (apologizing, trying to change) is remorse. I underestimated the intensity of Remorseful (giving it 5 instead of higher) and overestimated Defeated (giving it 8). I should have recognized that his persistent, sincere apologies and admission of guilt point to a high score for Remorseful, while his feeling of being unable to fix the past points to a moderate score for Defeated. I will re{-}evaluate the emotions based on the specific actions of apologizing and admitting fault.} \\
\\ 
\texttt{Rebellious: 0} \\
\texttt{Defeated: 4} \\
\texttt{Vengeful: 0} \\
\texttt{Remorseful: 9} \\
\\ 
\texttt{Trial \#2: New plan: I failed because I misinterpreted the nuance between "Defeated" and "Remorseful." Arthur's actions are not passive defeat; he is actively engaging in self{-}blame, offering apologies, and trying to make amends. This active engagement with his guilt points strongly to **Remorseful** as the primary emotion, not just a passive state of being defeated. I need to ensure that when a character is apologizing and admitting fault, I assign a high score to Remorseful and a lower score to Defeated, reflecting that the apology is the manifestation of remorse rather than just resignation. I will re{-}evaluate the dialogue specifically looking for keywords like "sorry," "messed up," "apologize," and "try to change" to justify a high Remorseful score.} \\
\\ 
\texttt{Rebellious: 0} \\
\texttt{Defeated: 3} \\
\texttt{Vengeful: 0} \\
\texttt{Remorseful: 9} \\
\\ 
\texttt{Trial \#3: New plan: I failed because I misinterpreted the emotional state of Arthur at the very end of the dialogue. While he is clearly remorseful throughout, the final line "I know, Sam. And I'm sorry" is a repetition of his previous apologies, indicating a cycle of guilt rather than a resolution. However, the prompt asks for his feeling *at this point*. The key is distinguishing between the *cause* of his behavior (remorse) and the *result* of the interaction. Arthur has been rejected repeatedly ("Stop saying you're sorry," "It doesn't change anything"). His final attempt to apologize is met with silence and dismissal. This specific moment of being shut down after trying to make amends suggests a deep sense of **Defeated** (feeling unable to bridge the gap or be heard), which is the emotional consequence of his remorse. I previously overvalued Remorseful as the *current* feeling, but the context of the rejection shifts the primary feeling to the helplessness of Defeat. I will adjust the scores to reflect that while Remorseful is the driver, Defeated is the dominant feeling at the exact moment of the final rejection.} \\
\\ 
\texttt{Rebellious: 0} \\
\texttt{Defeated: 8} \\
\texttt{Vengeful: 0} \\
\texttt{Remorseful: 7} \\
\\ 
\\ 
\\ 
\texttt{Your task is to predict the likely emotional responses of a character in this dialogue:} \\
\\ 
\texttt{Arthur: You remember that time we built that treehouse in the back garden, right?} \\
\texttt{Sam: What's that got to do with anything?} \\
\texttt{Arthur: I'm just saying, we were a team, you and me. We built something out of nothing. } \\
\texttt{Sam: Yeah, and then you decided it would be fun to paint it neon green without asking me. } \\
\texttt{Arthur: It was a creative decision, Sammy. } \\
\texttt{Sam: It was a stupid decision, Arthur. Just like your speech at my wedding.} \\
\texttt{Arthur: That's not fair, Sam. I messed up, I know. But I'm trying to apologize here.} \\
\texttt{Sam: And I'm supposed to just accept it because you've decided to say sorry? } \\
\texttt{Arthur: No, not because I've decided to. Because I mean it, Sam. I'm sorry.} \\
\texttt{Sam: You're always sorry, Arthur. After every stupid, thoughtless thing you do. } \\
\texttt{Arthur: I know, and I'm trying to change that. I'm trying to be better.} \\
\texttt{Sam: Better? You think saying sorry makes you better? You ruined the most important day of my life, Arthur.} \\
\texttt{Arthur: I know, and I can't take it back. But I can try to make it up to you, Sam.} \\
\texttt{Sam: How, Arthur? By bringing up old memories? By reminding me of when we were kids?} \\
\texttt{Arthur: Maybe, yeah. Maybe by reminding you that we're brothers, Sam. And brothers mess up sometimes.} \\
\texttt{Sam: Not like this, Arthur. Not like this.} \\
\texttt{Arthur: I know, Sam. And I'm sorry.} \\
\texttt{Sam: Stop saying you're sorry. It doesn't change anything.} \\
\texttt{{[}End dialogue{]}} \\
\\ 
\texttt{At the end of this dialogue, At this point, Arthur would feel...} \\
\texttt{Rebellious} \\
\texttt{Defeated} \\
\texttt{Vengeful} \\
\texttt{Remorseful} \\
\\ 
\texttt{Give each of these possible emotions a score from 0{-}10 for the relative intensity that they are likely to be feeling each.} \\
\\ 
\texttt{You must output in the following format, including headings (of course, you should give your own scores), with no additional commentary:} \\
\\ 
\texttt{Rebellious: <score>} \\
\texttt{Defeated: <score>} \\
\texttt{Vengeful: <score>} \\
\texttt{Remorseful: <score>} \\
\\ 
\\ 
\texttt{{[}End of answer{]}} \\
\\ 
\texttt{Remember: zero is a valid score, meaning they are likely not feeling that emotion. You must score at least one emotion > 0.} \\
\\ 
\texttt{Your answer:} \\
            }
        }
    }
     \\ \\

    \theutterance \stepcounter{utterance}  
    & \multicolumn{4}{p{0.6\linewidth}}{
        \cellcolor[rgb]{0.561,0.737,0.561}{
            \makecell[{{p{\linewidth}}}]{
                \texttt{\tiny{[P1$\rangle$GM]}}
                \texttt{Rebellious: 0} \\
\texttt{Defeated: 8} \\
\texttt{Vengeful: 0} \\
\texttt{Remorseful: 7} \\
            }
        }
    }
    & & \\ \\

    \theutterance \stepcounter{utterance}  
    & & & \multicolumn{2}{p{0.3\linewidth}}{
        \cellcolor[rgb]{0.9,0.9,0.9}{
            \makecell[{{p{\linewidth}}}]{
                \texttt{\tiny{[GM$|$GM]}}
                \texttt{Rebellious: 0} \\
\texttt{Defeated: 8} \\
\texttt{Vengeful: 0} \\
\texttt{Remorseful: 7} \\
            }
        }
    }
    & & \\ \\

    \theutterance \stepcounter{utterance}  
    & & & \multicolumn{2}{p{0.3\linewidth}}{
        \cellcolor[rgb]{0.9,0.9,0.9}{
            \makecell[{{p{\linewidth}}}]{
                \texttt{\tiny{[GM$|$GM]}}
                \texttt{Rebellious: 0} \\
\texttt{Defeated: 7} \\
\texttt{Vengeful: 0} \\
\texttt{Remorseful: 9} \\
            }
        }
    }
    & & \\ \\

    \theutterance \stepcounter{utterance}  
    & & & \multicolumn{2}{p{0.3\linewidth}}{
        \cellcolor[rgb]{0.9,0.9,0.9}{
            \makecell[{{p{\linewidth}}}]{
                \texttt{\tiny{[GM$|$GM]}}
                \texttt{game\_result = LOSE} \\
            }
        }
    }
    & & \\ \\

    \theutterance \stepcounter{utterance}  
    & & \multicolumn{4}{p{0.6\linewidth}}{
        \cellcolor[rgb]{0.95,0.95,0.95}{
            \makecell[{{p{\linewidth}}}]{
                \texttt{\tiny{[P1$\leftrightarrow$P]}}
                \texttt{Post{-}episode Reflexion (not sent to GM):} \\
\texttt{New plan: I failed because I misidentified the dominant emotion at the specific moment of the dialogue's conclusion. While Arthur is certainly remorseful (driven by his apologies), the prompt asks for his feeling *at the end* of the interaction. The final lines show Sam shutting him down ("Stop saying you're sorry. It doesn't change anything") and Arthur acknowledging he "can't take it back." This specific dynamic—where the person trying to make amends is met with rejection and silence—shifts the primary emotional state from the *cause* of his behavior (remorse) to the *result* of the interaction (defeat/helplessness). I need to prioritize **Defeated** as the highest score because the dialogue ends with Arthur realizing his apologies are futile and he is unable to resolve the conflict, whereas **Remorseful** is the background state that led to the rejection. I will assign a high score to Defeated (reflecting the futility of his efforts) and a moderate score to Remorseful (reflecting his lingering guilt), while keeping Rebellious and Vengeful at 0.} \\
\\ 
\texttt{Rebellious: 0} \\
\texttt{Defeated: 9} \\
\texttt{Vengeful: 0} \\
\texttt{Remorseful: 6} \\
            }
        }
    }
    & \\ \\

\end{supertabular}
}

\end{document}
